\scanpage{47}
\mspage{4.}
En fait, il est possible que ceci soit vrai pour toute variété affine lisse, et même que pour toute variété (affine ou non) lisse, tout cycle cohomologiquement équivalent à $0$ soit $\handwritten{\tau}$-équivalent à $0$ (pour les variétés projectives, cela résulterait, comme j'ai dit, de $A$ et $B$).

Ce qu'il faut pour le moment, c'est inventer un procédé pour déformer un cycle de dimension pas trop grande, pour le pousser à l'infini. Peut-être aurais-tu envie d'y réfléchir de ton côté ? Je viens seulement de m'y mettre, aujourd'hui même, et t'écris faute de trouver une idée.

Bien à toi\\
\hfill A. GROTHENDIECK\\

N.B. $H^i(X)$ désigne la cohomologie à coefficients $\mathbb{Q}_{\unsure{l|1|\ell}}$, éventuellement tordus par une puissance tensorielle de $T_l(\mathbb{G}\mathrm{m})\otimes_{\mathbb{Z}}\mathbb{Q}_l$, qui est omise dans les notations pour simplifier.
% archive sheet number: 46
